% Synthetic DCE - methods & results brief (formulas)
% Build:  pdflatex RESULTS.tex
\documentclass[11pt]{article}
\usepackage[margin=1in]{geometry}
\usepackage{amsmath,amssymb,booktabs,microtype,pifont,textcomp}
\usepackage[table]{xcolor}
\usepackage[colorlinks=true,linkcolor=blue!50!black,urlcolor=blue!50!black]{hyperref}

\newcommand{\yhat}{\hat{y}}
\newcommand{\roi}{\mathcal{M}}
\newcommand{\Var}{\operatorname{Var}}
\newcommand{\Ex}{\mathbb{E}}
\newcommand{\1}{\mathbb{1}}

\title{Synthetic DCE from Non-Contrast Prostate MRI\\\large Methods and Results Brief}
\author{Tier-1 static track --- UCSF cohort}
\date{6 August 2026}

\begin{document}
\maketitle

\begin{abstract}
\noindent
We predict a post-contrast DCE volume from three non-contrast sequences (T2w, DWI, ADC)
on 3{,}667 UCSF prostate exams. A 3D conditional GAN reaches within-gland co-localization
$r=0.482$, against $-0.043$ for the strongest trivial baseline (copying the T2w channel),
so the model predicts enhancement structure rather than passing anatomy through.
The substantive result is a decomposition: \textbf{within-gland structure is predictable;
per-patient enhancement amplitude is not}, and amplitude accounts for $51\%$ of the signal.
We show this is a property of the task rather than a tuning failure, and that the
unidentifiable component is orthogonal to the PI-RADS curve-type endpoint, which is
scale-invariant.
\end{abstract}

\section{Task}

Let $x=(x^{\text{T2w}},x^{\text{DWI}},x^{\text{ADC}})\in\mathbb{R}^{3\times D\times H\times W}$
be the stacked non-contrast input and $y\in[-1,1]^{1\times D\times H\times W}$ the target
DCE phase. We learn $G_\theta:x\mapsto\yhat$ with $\yhat=\tanh(\cdot)\in[-1,1]$.
$\roi_i\subset\Omega$ denotes the prostate mask of case $i$, $|\roi_i|\approx 5\%$ of the crop.

\section{Data and target selection}

\subsection{Cohort}

\begin{center}
\begin{tabular}{llr}
\toprule
Stage & Criterion & Remaining \\
\midrule
Patients with DCE & --- & 5{,}439 \\
Successfully staged & inputs resolvable, mask matches DCE grid & 4{,}975 \\
Enhancement QC (staging) & $\max_t \text{enh}(t) \ge 1.5$ & \textbf{4{,}246} \\
DWI b-value filter & DWI resolves to $b\ge600$ & 3{,}684 \\
Enhancement QC recorded & --- & \textbf{3{,}667} \\
\midrule
\multicolumn{2}{l}{Patient-level split} & 2805 / 312 / 550 \\
\bottomrule
\end{tabular}
\end{center}

Two cohort properties are findings rather than footnotes: $\approx12$--$13\%$ of patients
never enhance at any timepoint (failed or mistimed injection), and $13\%$ of DWI resolves
to $b{=}50$, which carries almost no diffusion weighting and therefore places a
\emph{different contrast} in the same input channel as the $b{=}1000$ the majority receive.

\subsection{Phase selection by elapsed time}

Enhancement relative to pre-contrast is
\begin{equation}
\text{enh}(t)\;=\;\frac{\Ex_{v\in\roi}\,[\,S_v(t)\,]}{\Ex_{v\in\roi}\,[\,S_v(0)\,]},
\end{equation}
measured as: flat at $\approx 1.0$ to $t\approx45$\,s; wash-in $55$--$90$\,s
($1.12\!\to\!2.00$); then a \textbf{plateau} of $2.1$--$2.5$ from $\approx 90$\,s past
$300$\,s. There is no sharp peak, so $\arg\max_t \text{enh}(t)$ lands on plateau noise
($\approx 280$\,s, into washout). Protocols are heterogeneous ($26$--$70$ phases,
$1.7$--$13$\,s cadence), so a fixed \emph{index} is a different physiological moment per
patient. We select by \emph{time},
\begin{equation}
j^\star \;=\; \arg\min_j \bigl| t_j - t_{\text{target}} \bigr|,
\qquad t_{\text{target}} = 120\,\text{s},
\end{equation}
which sits inside the plateau, past the timing-sensitive wash-in, at $\approx92\%$ of maximum.

\section{Preprocessing}

\paragraph{Geometry.} All modalities are resampled to a common physical grid
(\texttt{-{}-reference iso}) at spacing $(0.47,0.47,3.0)$\,mm, then center-cropped to
$32\times192\times192$, i.e.\ a fixed field of view of $90\times90\times96$\,mm for
\emph{every} patient. This matters: UCSF in-plane resolution spans $0.176$--$0.781$\,mm
across $256$--$1024$ matrices, so a fixed \emph{pixel} crop of 256 covered anywhere from
$40$\,mm to $320$\,mm of anatomy --- an $8\times$ range of scale, with 172 cases cropped
below the width of a prostate.

\paragraph{Intensity.} Per-image (single center, so no cross-scanner alignment).
For DCE, with body-tissue median $m$ and robust spread $s$,
\begin{equation}
y \;\longmapsto\; \operatorname{clip}\!\Bigl(\frac{y-m}{k\,s},\,-1,\,1\Bigr),\qquad k=1 .
\end{equation}
$k=1$ was selected by measuring within-prostate target dispersion: $k{=}3\Rightarrow 0.073$,
$k{=}2\Rightarrow0.109$, $k{=}1\Rightarrow\mathbf{0.218}$. Larger $k$ flattens the target
until ROI metrics become noise on a near-uniform field.

\section{Objective}

With $\phi$ a frozen 3D MedicalNet feature extractor,
\begin{equation}
\mathcal{L} \;=\; \lambda_1\,\mathcal{L}_{1}^{\text{ROI}}
\;+\;\lambda_s\bigl(1-\mathrm{SSIM}_{3D}(\yhat,y)\bigr)
\;+\;\lambda_s\bigl(1-\mathrm{SSIM}_{\roi}(\yhat,y)\bigr)
\;+\;\lambda_p\,\bigl\|\phi(\yhat)-\phi(y)\bigr\|_2^2 ,
\end{equation}
where the reconstruction term upweights the gland,
\begin{equation}
\mathcal{L}_{1}^{\text{ROI}}=\frac{\sum_{v\in\Omega} w_v\,\lvert \yhat_v-y_v\rvert}{\sum_{v\in\Omega} w_v},
\qquad
w_v = 1 + (\omega-1)\,\1[\,v\in\roi\,],\qquad \omega=10 .
\end{equation}
Without this the loss is dominated by background: the prostate is $\approx5\%$ of the crop.
For the latent diffusion models the objective is latent-space MSE, so ROI emphasis reaches
them only through the first-stage reconstruction.

\subsection{Flow matching}

On the straight-line (rectified) path with $t=0$ at the data and $t=1$ at the source,
\begin{equation}
z_t = \bigl(1-(1-\sigma_{\min})t\bigr)z_0 + t\,\epsilon,
\qquad
v^\star = \epsilon-(1-\sigma_{\min})z_0,
\qquad \epsilon\sim\mathcal{N}(0,I),
\end{equation}
with $\sigma_{\min}=0$ by default and $t\sim\mathcal{U}(0,1)$ sampled per example, so the
network learns the velocity field over the whole interpolation. The clean-latent identity
\begin{equation}
\hat z_0 \;=\; z_t - t\,v_\theta(z_t,t)
\end{equation}
holds exactly for any $\sigma_{\min}$ and is what the optional image-space anchor decodes
and supervises.

\section{Metrics}

All ROI statistics are computed \textbf{per case and then averaged}, never pooled over the
batch. Pooling mixes between-patient brightness into every statistic: on data where
per-patient brightness is predicted perfectly but within-gland structure is pure noise,
pooling reports $r=+0.977$ where the truth is $+0.020$.

\paragraph{Localization (faithfulness).} For case $i$ with $\bar y=\Ex_{v\in\roi_i}[y_v]$,
\begin{equation}
r_i=\frac{\sum_{v\in\roi_i}(\yhat_v-\bar{\yhat})(y_v-\bar y)}
{\sqrt{\sum_{v\in\roi_i}(\yhat_v-\bar{\yhat})^2}\;\sqrt{\sum_{v\in\roi_i}(y_v-\bar y)^2}},
\qquad
\texttt{roi\_pearson}=\frac{1}{N}\sum_{i=1}^{N} r_i .
\end{equation}

\paragraph{Texture and detail (realism).}
\begin{align}
\texttt{roi\_var\_ratio} &= \frac{1}{N}\sum_i \min\!\Bigl(\tfrac{\Var_{v\in\roi_i}[\yhat]}{\Var_{v\in\roi_i}[y]+\varepsilon},\,5\Bigr), \\[2pt]
\texttt{roi\_grad\_ratio} &= \frac{1}{N}\sum_i \min\!\Bigl(\tfrac{\Ex_{v\in\roi_i}\lVert\nabla\yhat\rVert}{\Ex_{v\in\roi_i}\lVert\nabla y\rVert+\varepsilon},\,5\Bigr),
\end{align}
both with target value $1$: below $1$ is over-smoothed, above $1$ is over-textured.
The intensity-histogram distance uses order statistics of the masked voxels,
\begin{equation}
\texttt{roi\_w1}=\frac{1}{2}\cdot\frac{1}{|\roi_i|}\sum_{k}\bigl|\yhat_{(k)}-y_{(k)}\bigr| ,
\end{equation}
the $\tfrac12$ mapping $[-1,1]\!\to\![0,1]$. These combine into a label-free proxy
\begin{equation}
\texttt{realism} = \operatorname{mean}\Bigl\{\,\nu(\texttt{var\_ratio}),\;\nu(\texttt{grad\_ratio}),\;1-\min(\texttt{w1},1)\,\Bigr\},
\qquad \nu(x)=1-\min\bigl(|x-1|,1\bigr).
\end{equation}

\paragraph{Fidelity.} $\mathrm{PSNR}=10\log_{10}\!\bigl(L^2/\mathrm{MSE}\bigr)$ with $L=2$
(the $[-1,1]$ range). Global SSIM/PSNR are $\approx99\%$ background and near-meaningless
here; the \emph{gap} between global and ROI values is the diagnostic.

\section{The central result: an identifiability decomposition}

\subsection{Baselines}

Three predictors that require no learning, with $\mu_i=\Ex_{v\in\roi_i}[y_v]$ and
$\bar\mu=\Ex_i[\mu_i]$:
\begin{equation}
\yhat^{\text{const}}_v=\bar\mu ,
\qquad
\yhat^{\text{level}}_v=\mu_i ,
\qquad
\yhat^{\text{t2w}}_v = x^{\text{T2w}}_v .
\end{equation}
\texttt{const} knows nothing; \texttt{level} is an \emph{oracle} that knows the correct
per-case brightness but has no internal structure; \texttt{t2w} is the identity baseline.

\begin{center}
\begin{tabular}{lccl}
\toprule
Predictor & $\mathrm{MAE}_{\roi}$ & \texttt{roi\_pearson} & Knows \\
\midrule
\texttt{const} & 0.2510 & --- & nothing \\
\texttt{t2w copy} & 0.3372 & $\mathbf{-0.043}$ & identity baseline \\
\textbf{MODEL (3D GAN)} & 0.2355 & $\mathbf{+0.482}$ & trained generator \\
\texttt{level} (oracle) & \textbf{0.1851} & --- & correct brightness, no structure \\
\bottomrule
\end{tabular}
\end{center}

\subsection{Variance decomposition}

Total ROI variance splits into a between-case (amplitude) and a within-case (structure) term:
\begin{equation}
\Var_{i,v}[\,y\,] \;=\; \underbrace{\Var_i[\mu_i]}_{\text{between-case, amplitude}}
\;+\; \underbrace{\Ex_i\bigl[\sigma_i^2\bigr]}_{\text{within-case, structure}},
\qquad \sigma_i^2=\Var_{v\in\roi_i}[y_v].
\end{equation}
Measured on the test set: $\operatorname{sd}_i[\mu_i]=0.231$ and $\Ex_i[\sigma_i]=0.226$, so
\begin{equation}
\frac{\Var_i[\mu_i]}{\Var_i[\mu_i]+\Ex_i[\sigma_i^2]}
=\frac{0.231^2}{0.231^2+0.226^2}
=\mathbf{0.51}.
\end{equation}

\subsection{Reading it}

\textbf{Structure is learned.} The T2w copy scores $\approx0$, so $r=0.482$ is genuine
localization and not laundered anatomy.

\textbf{Amplitude is not.} The model tracks per-patient brightness at only
$\operatorname{corr}_i(\mu_i,\hat\mu_i)=0.334$ and compresses between-case spread
$2.4\times$ ($0.231\to0.098$). Because amplitude is $51\%$ of the variance, both models lose
to the oracle-brightness baseline in MAE, and the flow model is worse than a \emph{constant}
predictor ($-0.018$) while still possessing real localization ($r=0.369$).

\textbf{Why this is the task, not the model.} Per-case amplitude depends on gadolinium dose,
injection rate, cardiac output, renal clearance and receiver gain --- none of which is
visible in T2w/DWI/ADC. We checked the DICOM headers for these variables: they were removed
by de-identification (dose, rate, patient weight, field strength absent; manufacturer
overwritten with \texttt{"NA"}). The information is not merely unused; it is not present.

\textbf{Why it may not matter.} PI-RADS curve types are scale-invariant: for any $\alpha>0$,
\begin{equation}
\text{type}\bigl(\alpha\,C(t)\bigr)=\text{type}\bigl(C(t)\bigr),
\end{equation}
since Type I (progressive), II (plateau) and III (washout) are defined by the \emph{shape}
of $C(t)$ after the peak. The unidentifiable component is orthogonal to the clinical endpoint.

\section{Results}

3D models, best checkpoint, held-out test ($n=550$).

\begin{center}
\begin{tabular}{lcc}
\toprule
Metric & 3D GAN & 3D flow (MedVAE) \\
\midrule
\texttt{roi\_pearson} $\uparrow$ & \textbf{0.482} & 0.369 \\
\texttt{ssim\_roi} $\uparrow$ & \textbf{0.397} & 0.276 \\
$\mathrm{MAE}_{\roi}$ $\downarrow$ & \textbf{0.236} & 0.269 \\
\texttt{roi\_var\_ratio} $\to1$ & \textbf{0.912} & 1.711 \\
\texttt{roi\_grad\_ratio} $\to1$ & 0.711 & \textbf{1.059} \\
\texttt{realism} $\uparrow$ & \textbf{0.845} & 0.714 \\
\texttt{roi\_pearson} PZ / TZ & 0.339 / 0.522 & 0.169 / 0.417 \\
\bottomrule
\end{tabular}
\end{center}

VAL $\approx$ TEST throughout. The GAN leads on faithfulness across four independent setups
(2D and 3D, pixel and latent). The flow's single advantage is texture
($\texttt{grad\_ratio}=1.06$ vs $0.71$): it is less smooth, which reads as more realistic
while localizing worse --- the faithfulness/realism tradeoff in one row.

\paragraph{Peripheral zone.} $r_{\text{PZ}}=0.339$ against $r_{\text{TZ}}=0.522$. PZ is
where most clinically significant cancer arises and where DCE drives the PI-RADS upgrade
rule, so this is the clinically weakest region and should be reported as such.

\section{Reliability}

\subsection{Checkerboard artifact (fixed)}

Transposed convolutions in all three decoders stamped a period-2 pattern into predictions.
Detector: the lag-1 autocorrelation of the first difference along a row,
\begin{equation}
\rho=\operatorname{corr}(d_k,\,d_{k+1}),\qquad d_k=y_{k+1}-y_k ,
\end{equation}
which tends to $-1$ under a pure checkerboard and to $-0.5$ for white noise.

\begin{center}
\begin{tabular}{lcc}
\toprule
& $\rho$ & HF power fraction \\
\midrule
Real DCE target & $+0.813$ & 0.0003 \\
GAN prediction (before) & $\mathbf{-0.598}$ & 0.0015 \\
GAN prediction (after) & $+0.278$ & 0.0003 \\
\bottomrule
\end{tabular}
\end{center}

The prediction alternated in sign at pixel scale \emph{more strongly than white noise}.
Removing it halved FID ($175.5\to88.9$), an independent confirmation.

\subsection{Run-to-run variance}

Two configuration-identical runs at the same seed differ. Training is not deterministic
(cuDNN autotuning; non-deterministic \texttt{atomicAdd} in 3D conv backward), adversarial
dynamics amplify the divergence, and best-checkpoint selection converts a bumpy validation
curve into a noisy discrete choice of epoch.

\begin{center}
\begin{tabular}{lc}
\toprule
Metric & $\Delta$ between identical runs \\
\midrule
\texttt{roi\_pearson} & 0.017 \\
\texttt{roi\_var\_ratio} & 0.055 \\
\texttt{p75\_corr} & \textbf{0.276} \\
\bottomrule
\end{tabular}
\end{center}

\paragraph{Rule.} Differences below $\approx0.02$ in \texttt{roi\_pearson} are unresolved at
$n=1$. \emph{Safe:} GAN $>$ flow ($0.113$, $\approx7\times$ noise); model $\gg$ t2w
($\approx30\times$); PZ $<$ TZ (paired within run). \emph{Not safe:} anything from
\texttt{p75\_corr} at $n=1$; the iso-vs-no-iso comparison ($0.020$, inside noise).

\paragraph{Why \texttt{p75\_corr} is the unstable one.} \texttt{roi\_pearson} averages $N=550$
per-case values, so variance is suppressed by $\sqrt{N}$. \texttt{p75\_corr} is a single
cross-patient correlation, and it measures the \emph{amplitude} dimension --- precisely the
component nothing in the loss constrains. \textbf{Reproducibility tracks identifiability}:
the well-determined component reproduces, the ill-determined one drifts.

\section{Tier-2: kinetics}

The multi-timepoint target is the enhancement curve $C_t(\cdot)$ per voxel. Unlike generic
video world models, the governing dynamics are \emph{known} --- the Tofts model,
\begin{equation}
\frac{dC_t}{dt} \;=\; K^{\text{trans}}\,C_p(t)\;-\;k_{ep}\,C_t(t),
\qquad k_{ep}=\frac{K^{\text{trans}}}{v_e},
\end{equation}
with $C_p$ the arterial input function. A continuous-time latent world model
(ODEWorld/PT-Flow) integrates a learned velocity field in \emph{physical} time,
\begin{equation}
\frac{dz_t}{dt}=v_\theta(z_t,t;z_0,c)
\quad\Longrightarrow\quad
z_T=z_0+\int_0^T v_\theta(z_t,t;z_0,c)\,dt ,
\end{equation}
where $z_0$ encodes the initial (pre-contrast) state and $c$ is a conditioning signal.
Our contribution would be to \emph{constrain} $v_\theta$ to the Tofts family with parameter
maps predicted from bpMRI, rather than learn it free-form --- a stronger inductive bias,
guaranteed physically plausible curves, and a direct answer to the limitation that
physics-informed ODE methods ``require explicit prior knowledge of the governing equations,''
which for DCE we have.

\paragraph{Goal state.} Acquisition duration (measured, $n=1367$): $p_5=222$\,s,
$p_{50}=300$\,s, $p_{95}=380$\,s. Coverage is $99.0\%$ at $180$\,s, $\mathbf{91.9\%}$ at
$240$\,s, and only $49.7\%$ at $300$\,s. We therefore anchor at $t=240$\,s: past plateau
onset by $\approx150$\,s, inside the window where Types II and III separate, while retaining
$92\%$ of the cohort.

\section{Paper partitions and venue paths}

The work separates into three claims with different evidence requirements, different
readiness, and different venues. They are not sequential --- \textcircled{2} is nearly
complete on existing data while \textcircled{1} waits on labels.

\begin{center}
\renewcommand{\arraystretch}{1.25}
\begin{tabular}{p{0.9cm}p{4.4cm}p{4.6cm}p{3.5cm}}
\toprule
& \textbf{Claim} & \textbf{Still needs} & \textbf{Venue} \\
\midrule
\textcircled{1} &
Single-phase synthetic DCE is clinically usable; multi-vendor, PI-RADS endpoint &
Hanxue's labels; reader study; multi-seed repeats &
Radiology / Lancet Dig.\ Health / Insights into Imaging \\
\textcircled{2} &
Standard synthesis metrics do not measure case-specific fidelity in ill-posed conditional
generation &
A within-model sweep; $\ge3$ seeds; a \emph{second domain} &
ICLR 2027 (tight) $\to$ ICML / NeurIPS \\
\textcircled{3} &
A physics-constrained velocity field beats a free-form one for contrast kinetics &
Multi-phase staging (done); Tofts parameterization; interpolation ablation &
NeurIPS world-models workshop $\to$ MICCAI 2027 \\
\bottomrule
\end{tabular}
\end{center}

\paragraph{\textcircled{1} --- clinical.} The models exist and the pipeline is audited; the
blocker is entirely the endpoint. The overall PI-RADS category is essentially absent from the
report export ($92/10{,}538$ once a boilerplate footer URL is excluded), but the PI-RADS v2.1
\emph{DCE component} is binary positive/negative and appears in $\approx6{,}400$ reports, of
which $\approx2{,}374$ fall in our QC-passing cohort. That is the natural endpoint --- it is
the one component synthetic DCE is responsible for, and $\kappa$ applies directly to a binary
call. We are waiting on Hanxue's labels rather than mining the free text ourselves.

\paragraph{\textcircled{2} --- identifiability and evaluation.} The strongest near-term
result, and it is \emph{label-free}: the amplitude/structure decomposition (\S5), the finding
that both models lose to an oracle-brightness baseline while retaining real localization, and
the observation that \textbf{reproducibility tracks identifiability}. The framing must avoid
the perception--distortion tradeoff (Blau \& Michaeli, CVPR 2018), which is established ---
the contribution is a \emph{diagnostic for whether a conditional model used its conditioning
at all}, together with evidence that the reported metrics cannot distinguish that case.
A single medical domain will read as an application paper, so a second domain is the
gating requirement, not a nicety.

\paragraph{\textcircled{3} --- kinetics.} The novelty is not porting a world model to a new
organ (that is a MICCAI paper, and the source method's authors would rightly be co-cited
rather than superseded); it is that \emph{we know the governing equation}. ODEWorld's own
related work identifies the requirement for ``explicit prior knowledge of the true governing
equations'' as the limitation of physics-informed ODE methods --- for DCE that prior knowledge
is the Tofts model. The experiment is therefore free-form $v_\theta$ versus Tofts-constrained
$v_\theta$ versus Tofts-plus-residual, under irregular temporal sampling.

\paragraph{Timing.} Approximate and worth confirming: ICLR 2027 closes around late September
2026 ($\approx7$ weeks out), NeurIPS workshops around September--October, MICCAI 2027 around
February 2027 ($\approx6$ months). Only \textcircled{2} is plausibly ICLR-ready in that
window, and only if a second domain moves quickly; \textcircled{3} fits the workshop now and
a full venue in the MICCAI cycle. A rushed ICLR rejection would also consume the MICCAI
window, so the go/no-go should be taken on the pre-contrast conditioning result rather than
on ambition.

\paragraph{On the Tier-1 $\to$ Tier-2 gate.} The project's onboarding sets a hard gate:
do not launch Tier 2 until PI-RADS $\kappa\ge0.6$. That gate is currently \emph{unmeasurable}
rather than failed, since the overall category is absent from the data. Two developments
argue for revisiting it: the scarcity rationale behind it has evaporated (UCSF supplies
$\approx8\times$ the planned Tier-2 multi-phase pool at $\approx10\times$ the temporal
density), and curve-type classification is scale-invariant, so the one component we cannot
predict is orthogonal to the Tier-2 endpoint. A defensible revision is to gate on
\emph{DCE-positivity} $\kappa$ instead of overall PI-RADS $\kappa$.

\section{Anticipated technical questions}

\paragraph{Q1. Is \texttt{roi\_pearson} $=0.482$ inflated by \emph{zonal} contrast rather
than focal enhancement? A model that only learns ``TZ is brighter than PZ'' would score well
without detecting anything.}
This is the sharpest objection and the data partly answers it. If the whole-gland correlation
were driven by the PZ/TZ contrast, we would expect it to greatly exceed the within-zone
values, since the between-zone term would be doing the work. It does not:
$r_{\text{gland}}=0.482$ sits \emph{between} $r_{\text{PZ}}=0.339$ and $r_{\text{TZ}}=0.522$,
and both within-zone correlations are solidly positive. So there is genuine within-zone
structure, not merely zonal architecture. What this does \emph{not} establish is
lesion-level detection --- reproducing normal zonal anatomy correctly still scores well.
That requires lesion labels and a lesion-level readout, which we do not yet have.

\paragraph{Q2. You report a mean over 550 cases with no interval. What is the spread?}
Not yet reported, and it should be. \texttt{roi\_pearson} is $\tfrac1N\sum_i r_i$; the
per-case $r_i$ distribution and a bootstrap CI over cases are computed from data we already
have. Until then, the run-to-run figure ($\pm0.017$) bounds \emph{run} variability but says
nothing about \emph{case} variability, and the two should not be conflated.

\paragraph{Q3. Why Pearson? The signal--concentration relationship is nonlinear.}
Correct, and it is a real limitation. MR signal relates to gadolinium concentration through
$T_1$ relaxation and the spoiled-gradient-echo signal equation, so a linear correlation
measures the wrong thing if the nonlinearity is severe within the ROI dynamic range. Spearman
$\rho$ would be the rank-based alternative and is trivial to add. We report Pearson because
it is what the synthesis literature uses; reporting both would be more defensible.

\paragraph{Q4. Does the $51\%$ decomposition assume independence between the two terms?}
No. The law of total variance,
$\Var[Y]=\Ex_i\!\bigl[\Var[Y\mid i]\bigr]+\Var_i\!\bigl[\Ex[Y\mid i]\bigr]$, is an identity;
no independence is required. One precision caveat: we substituted
$\bigl(\Ex_i[\sigma_i]\bigr)^2=0.226^2$ for $\Ex_i[\sigma_i^2]$. By Jensen
$\Ex[\sigma^2]\ge(\Ex[\sigma])^2$, so the true within-case term is at least that large and
\textbf{$51\%$ is an upper bound on the amplitude share}. The qualitative conclusion --- that
amplitude is roughly half the signal --- is unaffected.

\paragraph{Q5. Best checkpoints are selected on \texttt{ssim\_roi}, which is smoothness-biased.
Is ``the GAN wins'' partly a selection artifact?}
Partly, and this should be stated. \texttt{ssim\_roi} rewards smooth predictions, and the GAN
is the smoother model ($\texttt{grad\_ratio}=0.71$ vs $1.06$). Selecting on
\texttt{realism} or \texttt{balanced} would plausibly favour the flow. Two mitigations: the
GAN also leads on \texttt{roi\_pearson}, which is not smoothness-biased, and the margin
($0.113$) is $\approx7\times$ the run-to-run noise. But a selection-metric ablation
(same runs, re-selected under each criterion) is the clean answer and has not been done.

\paragraph{Q6. The LDM's ROI weighting only enters through the first stage. Is the diffusion
objective effectively unweighted?}
Yes, and it is a genuine asymmetry. The GAN's reconstruction term is ROI-weighted at
$\omega=10$; the LDM trains on latent-space MSE with no clean mapping back to the mask, so
ROI emphasis reaches it only via the frozen autoencoder's reconstruction. Combined with a
MedVAE first stage never trained on prostate DCE, the flow is handicapped in ways not
isolated from its objective. ``The GAN is better'' is therefore a statement about these two
configurations, not about GANs versus flows in general.

\paragraph{Q7. Were the 3D runs converged? 40 epochs versus 100 for 2D looks arbitrary.}
It was a wall-clock choice ($\approx670$\,s/epoch $\times$ 40 $\approx$ 7.4\,h), not a
convergence criterion, and convergence was not demonstrated. Against that: the GAN
\emph{degrades} late (discriminator wins, $d\to0.03$, $\mathrm{adv}\to4.4$), and the selected
best checkpoint clearly beats the final epoch ($0.482$ vs $0.419$), so longer training is not
obviously the missing ingredient. The \texttt{val\_trajectory.jsonl} now logged per run will
answer this directly.

\paragraph{Q8. Are the $12$--$13\%$ ``non-enhancers'' biology, or registration failure?}
Not yet separated, and it is a real confound. Enhancement QC is computed as a mask-mean ratio,
so a mis-registered mask would depress it exactly like a failed injection. We now capture a
per-case registration quality (\texttt{reg} in \texttt{dce\_times.json}, cohort median $0.459$,
$p_{90}=1.694$) but have not yet correlated it against \texttt{enh\_max}. If the
non-enhancers concentrate in the poor-registration tail, the ``failed injection'' interpretation
is wrong and the exclusion criterion needs rethinking.

\paragraph{Q9. FID on 550 grayscale MR volumes with an ImageNet Inception backbone --- is that
meaningful?}
Weakly. FID is biased at small sample sizes (KID is the low-$n$ alternative), and Inception
features are natural-image features. A RadImageNet or medical-SSL backbone would be the
defensible choice. We use FID because it is the field standard and because it responded
sharply and correctly to a known artifact (the checkerboard: $175\to89$), which is evidence
it is sensitive to \emph{something} real. It should not carry a headline claim.

\paragraph{Q10. \texttt{realism\_score} is an unvalidated composite with equal weights.}
Agreed --- it is a ranking heuristic for choosing between checkpoints without a reader
session, not a validated metric, and no weighting study was done. It should not appear in a
results table without that caveat; the components (\texttt{var\_ratio}, \texttt{grad\_ratio},
\texttt{w1}) should be reported individually.

\paragraph{Q11. \texttt{var\_ratio} and \texttt{grad\_ratio} are capped at 5. Doesn't that
bias the mean?}
Yes, by construction. The cap exists because a near-flat target ROI puts a small denominator
in a ratio, and uncapped cases reached the thousands and destroyed the mean. The estimator is
therefore biased toward 5 from above and should be reported with the median and the fraction
of cases hitting the cap --- neither of which we currently log.

\paragraph{Q12. Your target is normalized signal, not gadolinium concentration. How does that
square with a Tofts-based Tier-2?}
It does not, cleanly. Tofts governs $C_t(t)$, and converting signal to concentration requires
a baseline $T_1$ map (e.g.\ variable flip angle) that this cohort does not include. The
workable fallback is relative enhancement, which makes the recovered parameters
\emph{semi-quantitative} --- defensible if stated plainly, not if presented as $K^{\text{trans}}$
in physical units. An arterial input function is a second unresolved requirement, and the
$90$\,mm crop may not contain a usable artery.

\paragraph{Q13. Why not force determinism instead of measuring the noise?}
\texttt{torch.use\_deterministic\_algorithms(True)} would raise on several 3D convolution
kernels that have no deterministic implementation, and where it works it is substantially
slower. More importantly it would suppress a real signal: the divergence is concentrated in
the amplitude dimension precisely because nothing constrains it, so the instability is
\emph{informative}. The right response is $\ge3$ seeds with mean $\pm$ sd, not determinism.

\paragraph{Q14. A random $15\%$ patient split of a single center --- is the effective sample
size really 550?}
Probably smaller. Protocols cluster (\texttt{DCE PROSTATE}, Siemens TWIST, GE DISCO), and the
split is not stratified by protocol, so test cases are not fully independent draws.
Stratifying the split by \texttt{series} --- now captured in \texttt{stage\_meta.json} --- and
reporting per-protocol metrics would be the fix, and would double as the multi-vendor
generalization result the project promises.

\paragraph{Q15. How were the prostate masks produced, and were they checked?}
They are automatic segmentations delivered with the cohort (\texttt{pilot\_log.json} records
the strategy and resulting volume per case); we have not manually verified any of them.
Every ROI metric in this document is conditioned on those masks being right, so mask error
propagates directly into the headline numbers.

\section{Caveats}

2D and 3D ROI metrics are not comparable (the 2D dataset keeps only gland-containing slices,
so the denominators differ). FID was measured on the last checkpoint, not the best.
Every result predating the 2026-07 audit is superseded. No claim is made that additional
data would close the amplitude gap: an earlier conclusion to that effect was measured on a
configuration with known defects and has been retracted.

\end{document}
